\section{Security design}\label{sec:security}

This section describes the controls as designed and built. The companion \emph{Security self-assessment report}
evaluates them, rates the remaining risks and lists what is out of scope.

\subsection{Threat model}

\begin{table}[htbp]
  \centering
  \small
  \begin{tabularx}{\linewidth}{@{}p{26mm}XX@{}}
    \toprule
    Asset & Threat & Primary controls \\
    \midrule
    Sandbox money & The model, or text injected through merchant names or chat, moves money & Tiers and the gate; signed, single-use, expiring authorizations only the user can create; terms re-read from storage (\cref{sec:confirmation}) \\
    Truth of advice & The model invents or bends numbers, or hides a caveat & Every number must be a fact, every warning mentioned; otherwise deterministic narration (\cref{sec:verifier}) \\
    Suitability & A plan more aggressive than the finances allow & Permitted profile $=$ min(tolerance, capacity), recomputed server-side; hard rules H1--H5 \\
    Evidence & Audit events edited, removed or forged & Hash chain; erasure only through committed tombstones (\cref{sec:chain}) \\
    Identity & A caller claims another user & Cookie sessions resolved server-side; identity header rejected; Origin check (\cref{sec:http}) \\
    Personal data in statements & Stored, returned, logged or sent to a provider & Redaction at import (ADR-028) \\
    Data at rest & A copy of the database, its WAL, a backup or the volume & Every personal payload sealed; key kept off the volume (\cref{sec:at-rest}) \\
    Real data on a sandbox & A visitor uploads a real statement (open sign-up) & Sandbox notice and acknowledgement; sealing; deletion (\cref{sec:notice}) \\
    Third-party text & Instructions in a receipt, a concept card, a Fintual goal name or a connected service's answer steer the model & Wrapped as untrusted text; no tool writes; numbers are facts; the narration never repeats goal names (S24, S30, S32, S36) \\
    Provider credentials & A Fintual password or token, an MCP token or a Fintoc link leaks, or the agent links an account & Password used once and dropped; tokens sealed and bound to user and adapter; never in a response, log, prompt or audit; link and unlink are T3 (S31, S38) \\
    Connected servers & A user-added address reaches the internal network (SSRF), or a ``read-only'' tool writes & HTTPS and public addresses only, no redirects, size and time limits; tools enabled only if read-only, argument-free and without a write verb; only the user's sync calls them (S35, S37) \\
    \bottomrule
  \end{tabularx}
  \caption{Assets, threats and the controls that answer them (\texttt{docs/SECURITY.md}).}\label{tab:threats}
\end{table}

\subsection{Trust boundaries}

\Cref{fig:trust} shows the public demo: a browser reaches Cloudflare's edge over TLS, a tunnel opened from the host
carries the request to the web container, and nginx proxies \code{/api} to the API. Six boundaries matter.
\textbf{B1} separates untrusted visitors (and their files) from the edge; \textbf{B2} is the tunnel into the host;
\textbf{B3} is the API's own perimeter, where the identity guard, the rate limits and input validation run; \textbf{B4}
is the data key's separate mount; \textbf{B5} is the path to model providers, crossed only in live mode and only by
redacted text, and whose answers are untrusted; \textbf{B6} is the host itself, outside the scope of the application.

\begin{figure}[htbp]
  \centering
  % Trust boundaries of the public demo (deployment files on ops/deploy-neural-ledger, data key from this branch).
\begin{tikzpicture}[x=1mm,y=1mm]
  \tikzset{tb/.style={num,fill=nlfilldark}}
  \node[ext,text width=24mm] (user) at (0,0) {\textbf{Browser}\\untrusted: text, files, headers};
  \node[box,text width=25mm] (cf) at (33,0) {\textbf{Cloudflare edge}\\TLS; rate rule on auth POSTs};
  \node[box,text width=24mm] (cfd) at (69,0) {cloudflared\\outbound tunnel};
  \node[box,text width=24mm] (web) at (69,-24) {web (nginx)\\headers, \code{/api} proxy};
  \node[core,text width=38mm,align=left] (api) at (117,-24) {\textbf{api}\\identity guard, rate limits\\permission gate, verifier\\keys held in memory};
  \node[store] (vol) at (104,-50) {data volume\\SQLite + WAL};
  \node[keybox] (sec) at (130,-50) {secret\\file};
  \node[ext,text width=40mm] (prov) at (116,14) {\textbf{Model providers} (live mode)\\replies are untrusted: checked by the gate and the verifier};
  \begin{scope}[on background layer]
    \node[zone,fit=(cfd)(web)(api)(vol)(sec),inner sep=5mm] (host) {};
    \node[boundary,fill=white,fit=(api)(vol)(sec),inner sep=2mm] (cont) {};
  \end{scope}
  \node[anchor=north west,font=\footnotesize\bfseries] at (host.north west) {Host (nixbox), Docker};
  \node[anchor=north east,font=\scriptsize] at (cont.south east) {api container};
  \draw[flow] (user) -- (cf);
  \draw[flow] (cf) -- (cfd);
  \draw[flow] (cfd) -- (web);
  \draw[flow] (web) -- (api);
  \draw[flow] (api.south -| vol.north) -- (vol.north);
  \draw[flow] (sec.north) -- (sec.north |- api.south);
  \draw[flow,dashed] (api.north) -- node[label,right]{HTTPS, redacted text} (prov.south);
  \node[tb] at ($(user.east)!0.5!(cf.west)+(0,5)$) {B1};
  \node[tb] at ($(cf.east)!0.5!(cfd.west)+(0,5)$) {B2};
  \node[tb] at ($(web.east)!0.5!(api.west)+(0,5)$) {B3};
  \node[tb] at ($(sec.north)+(5,6)$) {B4};
  \node[tb] at ($(api.north)+(-5,5)$) {B5};
  \node[tb] at ($(host.south west)+(5,5)$) {B6};
  \node[anchor=west,font=\scriptsize] at ($(host.south west)+(9,5)$) {root on the host};
\end{tikzpicture}

  \caption{Trust boundaries of the public demo. The deployment override that removes host ports and adds the tunnel is on \texttt{ops/deploy-neural-ledger}.}\label{fig:trust}
\end{figure}

\subsection{Agent permissions}

The gate (\code{domain/permissions/policy.py}) knows four tiers: T0 reads, T1 analyzes and proposes, T2 executes an
approved sandbox action, T3 is prohibited and always blocked and audited (\cref{tab:tools}). The policy
\code{pp-0.2} allows two actions, each with a fixed route: a savings rule from \code{checking} to
\code{sandbox_savings} (at most 1{,}000{,}000 per action) and a liquidity floor that stays on \code{checking} (at most
100{,}000{,}000). The simulated bank checks the destination again on its side. Text from statements, merchants and
the profile's display name reaches the model only wrapped as \texttt{\{"untrusted\_text": \dots\}}, and profile writes are T3
for the agent, so injected text cannot persist as instructions. The newer sources of third-party text follow the
same rule. Receipts, concept-card passages (at most 1{,}500 characters), Fintual goal names and connected services'
answers (at most 4{,}000 characters) reach the model only as untrusted text. No tool lets the agent store,
correct, link, sync or delete any of them.

\subsection{Accounts, sessions and the sandbox notice}\label{sec:notice}

Passwords are hashed with scrypt ($N = 2^{14}$, $r = 8$, $p = 1$, 16-byte salt; \code{adapters/outbound/system/password.py});
stored parameters are bounded, so a tampered row cannot ask for unbounded memory. A sign-in with an unknown email does the
same scrypt work on a dummy hash, and both failures return byte-identical bodies (S20); sign-up, however, answers 409
for a registered email. Sessions are described in \cref{sec:http}.

Sign-up is open (option B). The sign-up page always shows a notice that cannot be dismissed (``This is a hackathon
sandbox... do not upload real bank statements and do not enter real personal data'') and a required checkbox. The
server enforces it: \code{SignupRequest.sandbox_acknowledged} must be \code{true}, or the answer is 422 on that field and
nothing is stored; \code{USER_REGISTERED} records which notice was accepted (\code{sandbox_notice=sandbox-1}). The
statement upload screen repeats the warning above the file picker.

Sign-up also checks the email's domain (\code{NL_EMAIL_CHECKS}): it refuses a disposable-inbox domain (the maintained
\code{disposable-email-domains} list, offline, parent domains included) and, with \code{dns}, a domain that does not
exist or receives no mail. Only the domain is sent to the resolver; a DNS timeout lets the email through.

\subsection{Data at rest}\label{sec:at-rest}

The SQLite adapter seals every payload listed in \cref{tab:schema} with AES-256-GCM (\code{FieldCipher} in
\code{adapters/outbound/persistence/field_cipher.py}): a random 12-byte nonce per value, stored as
\code{nlenc2:} followed by base64 of nonce, ciphertext and tag. The associated data is \texttt{<table>|<user\_id>|<row id>},
so a value copied into another user's row, or into another table, fails to open instead of being read. A value that
is not sealed is refused when read: after the migration, plain text in the file can only be an edit made around the
application, so it cannot downgrade a field to clear text. The email, which must stay searchable, is replaced by
\code{email_index} $=$ HMAC-SHA256 of the normalized address; the address itself is only inside the sealed account
payload. The audit record is sealed whole, its hashes included, so the file holds no hash against which a guessed
event could be tested.

\subsubsection{Key hierarchy}\label{sec:keys}

\code{NL_DATA_KEY} is 32 random bytes. It is never used directly: HKDF-SHA256 derives one 256-bit subkey per purpose
(\cref{fig:keys}). Values sealed before option B with the raw key (format \code{nlenc1:}) still open, and the migration
rewrites them. Two other secrets are separate roots: \code{NL_HMAC_SECRET} signs authorizations, and session tokens and
passwords are only ever stored as digests.

\begin{figure}[htbp]
  \centering
  % Key hierarchy (bootstrap/data_key.py, adapters/outbound/persistence/field_cipher.py) and the other secrets.
\begin{tikzpicture}[x=1mm,y=1mm]
  \tikzset{src/.style={box,text width=33mm,align=left,minimum height=8mm}}
  \node[font=\footnotesize\bfseries,anchor=south west] at (-19,21) {key sources, in this order};
  \node[src] (s1) at (0,15) {\textbf{1.} \code{NL_DATA_KEY_FILE}\\\scriptsize if set, it wins; must exist};
  \node[src] (s2) at (0,3) {\textbf{2.} Docker secret\\\scriptsize\code{/run/secrets/nl_data_key}};
  \node[src] (s3) at (0,-9) {\textbf{3.} \code{NL_DATA_KEY}\\\scriptsize variable};
  \node[src,dashed] (s4) at (0,-24) {\textbf{local runs only:} \code{<db>.key} on the data volume, with a warning};
  \node[keybox,text width=27mm] (master) at (42,-6) {\code{NL_DATA_KEY}\\32 random bytes\\(64 hex characters)};
  \foreach \s in {s1,s2,s3} \draw[flow] (\s.east) -- (master.west);
  \draw[flow,dashed] (s4.east) -- (master.south west);
  \node[box,text width=33mm] (hk1) at (82,8) {HKDF-SHA256, info\\\scriptsize\code{neural-ledger/v2/field-encryption}};
  \node[box,text width=33mm] (hk2) at (82,-20) {HKDF-SHA256, info\\\scriptsize\code{neural-ledger/v1/blind-index/email}};
  \draw[flow] (master.east) -- (hk1.west);
  \draw[flow] (master.east) -- (hk2.west);
  \node[keybox,text width=21mm] (fk) at (124,8) {field key\\256 bit};
  \node[keybox,text width=21mm] (ik) at (124,-20) {email index\\key, 256 bit};
  \draw[flow] (hk1) -- (fk);
  \draw[flow] (hk2) -- (ik);
  \node[core,text width=24mm] (aes) at (124,27) {AES-256-GCM\\seal, open\\\code{nlenc2:}};
  \node[core,text width=24mm] (hmac) at (124,-39) {HMAC-SHA256\\of the email:\\\code{email_index}};
  \draw[flow] (fk) -- (aes);
  \draw[flow] (ik) -- (hmac);
  \draw[weak] (master.north) |- node[label,pos=0.72]{opens legacy \code{nlenc1:} values only} (aes.west);
  \node[keybox,text width=33mm,align=left] (hmacsec) at (0,-52) {\code{NL_HMAC_SECRET}\\variable, at least 16 characters};
  \node[core,text width=40mm] (sign) at (50,-52) {HMAC-SHA256 of the canonical\\action terms: \code{action_hash}};
  \draw[flow] (hmacsec) -- (sign);
  \node[box,text width=42mm,align=left] (other) at (110,-56) {Not keys: the session token is stored as its SHA-256 digest; the password as scrypt ($N{=}2^{14}$, $r{=}8$, $p{=}1$).};
\end{tikzpicture}

  \caption{Key hierarchy. Thick-ruled boxes are secret material; shaded boxes are what the keys are used for.}\label{fig:keys}
\end{figure}

\Cref{alg:keys} is how the API finds the key and what a deployment refuses (\code{bootstrap/data_key.py},
\code{_data_cipher} in \code{bootstrap/container.py}). A refusal happens before the database is opened, so it leaves no
file behind. Once a key is chosen, \code{verify_data_key} opens a sealed known value in the \code{meta} table (or, in a
file that has none, the first sealed value of each column) and refuses a key that does not open it.

\begin{algorithm}[htbp]
  \caption{Choosing the data key (\texttt{resolve} and \texttt{check}).}\label{alg:keys}
  \KwIn{settings $S$, \code{/proc/self/mountinfo}}
  \KwOut{a field cipher, or a refusal}
  $f \gets S.\mathit{data\_key\_file}$ or \code{/run/secrets/nl_data_key}\;
  \uIf{$f$ is a file}{$k \gets$ its content}
  \uElseIf{$S.\mathit{data\_key\_file}$ was set}{refuse ``not a file''}
  \uElseIf{\code{NL_DATA_KEY} is set}{$k \gets$ \code{NL_DATA_KEY}}
  \lElse{$k \gets$ none}
  \If{$S.\mathit{deployment}$ is deployed}{
    \lIf{$k$ is none}{refuse ``no data key''}
    \lIf{$k$ came from a file on the database's mount (or in its directory)}{refuse ``key on the data volume''}
    \lIf{\code{<db>.key} still exists}{refuse ``run migrate-data, then delete it''}
    \lIf{\code{NL_HMAC_SECRET} is the development default}{refuse}
  }
  \lIf{$k$ is not none}{\Return cipher($k$)}
  \lIf{live mode}{refuse ``live mode needs NL\_DATA\_KEY''}
  warn loudly; \Return cipher(key file \code{<db>.key}, created once) \Comment*[r]{local runs only}
\end{algorithm}

\subsubsection{Migration and key rotation}

\code{neural-ledger migrate-data} (\code{bootstrap/migrate.py}, \code{adapters/outbound/persistence/sqlite_migration.py})
moves a database to the configured key once (\cref{alg:migrate}). It runs before the application is built, so the
refusals above do not stop it; the operator stops the API first. Pointing \code{--old-key-file} at the previous key
rotates keys the same way.

\begin{algorithm}[htbp]
  \caption{Migration (\texttt{migrate}).}\label{alg:migrate}
  \KwIn{database $D$, new cipher $c'$, old ciphers $O$ (from \code{<db>.key} or \code{--old-key-file}), dry run flag}
  \KwOut{counts per table: current, re-keyed, sealed from plain text}
  \textsc{begin immediate}\;
  \If{\code{PRAGMA user_version} $< 1$}{rebuild \code{user_accounts} (email $\to$ \code{email_index}), and \code{proposals}, \code{authorizations}, \code{audit_events} (add \code{user_id})}
  \ForEach{sealed column $(t, \mathit{aad})$ and row $r$}{
    $v \gets r.\mathit{value}$\;
    \uIf{$v$ is not sealed}{$x \gets v$ \Comment*[r]{plain text from before}}
    \uElseIf{$c'$ opens $v$}{\lIf{$v$ is \code{nlenc2:}}{count current; \KwContinue} $x \gets c'.\mathrm{open}(v)$}
    \uElseIf{some $c \in O$ opens $v$}{$x \gets c.\mathrm{open}(v)$}
    \lElse{\textsc{rollback}; refuse ``opens with no given key''}
    write $c'.\mathrm{seal}(x, \mathit{aad}(r))$\;
  }
  write the key check sealed with $c'$\;
  \eIf{dry run}{\textsc{rollback}}{\textsc{commit}; \textsc{vacuum}; WAL checkpoint (truncate)}
\end{algorithm}

The transaction makes a failure or a dry run leave the file untouched; running the command again changes nothing; and
VACUUM rewrites the file, so free pages keep no copy of the old plain text or old ciphertext. It ends by telling the
operator to delete the old key file, which deployed mode refuses to start next to.

\subsection{Account deletion}\label{sec:deletion}

\code{DELETE /v1/account} takes the password again and a typed confirmation, \code{DELETE} or the account's email
(\cref{fig:deletion}). It deletes every row the user owns, in one transaction, and turns the user's audit events into
tombstones that one \code{ACCOUNT_DELETED} event commits to (\cref{alg:erase}, \cref{sec:chain}). Every session ends
because its row is gone, and the response clears the cookie. Conversations and snapshots, which live only in memory,
are dropped too. The event that records the deletion has no user id and no session, so nothing in the database points
at the user any more.

\begin{figure}[htbp]
  \centering
  % Account deletion as a sequence (adapters/inbound/http/app.py, application/use_cases/accounts.py, persistence/sqlite.py).
\begin{tikzpicture}[x=0.985mm,y=1mm]
  \tikzset{actor/.style={box,text width=23.5mm,minimum height=9mm},
           msg/.style={->,thick},ret/.style={->,densely dashed},
           note/.style={draw,fill=nlfill,align=left,inner sep=3pt,font=\scriptsize}}
  \foreach \x/\n/\t in {0/ui/{User\\(profile page)}, 34/http/{HTTP API\\guard, rate limit}, 68/uc/{Account\\UseCases}, 102/er/{AccountErasure\\one transaction}, 136/au/{Audit chain}} {
    \node[actor] (\n) at (\x,0) {\t};
    \draw[densely dotted] (\n.south) -- (\x,-126);
  }
  \draw[msg] (0,-12) -- node[label,above]{\code{DELETE /v1/account}} (34,-12);
  \node[font=\scriptsize,anchor=north] at (17,-13) {confirm, password};
  \draw[msg] (34,-22) -- node[label,above]{\code{delete_account}} (68,-22);
  \node[note,anchor=north west,text width=29mm] at (69,-25) {confirm is DELETE or the email; scrypt check of the password};
  \draw[msg] (68,-44) -- node[label,above]{\code{erase(keep_account=False)}} (102,-44);
  \node[note,anchor=north west,text width=29mm] at (103,-47) {BEGIN IMMEDIATE; delete the user's rows in 17 tables (\code{secure_delete})};
  \draw[msg] (102,-66) -- node[label,above]{tombstone A's events} (136,-66);
  \draw[msg] (102,-75) -- node[label,above]{append \code{ACCOUNT_DELETED}} (136,-75);
  \node[note,anchor=north west,text width=29mm] at (103,-78) {detail: count and SHA-256 of the erased sequences, row counts; no user id};
  \node[note,anchor=north west,text width=29mm] at (103,-95) {COMMIT; WAL checkpoint};
  \draw[ret] (102,-104) -- node[label,above]{row counts} (68,-104);
  \draw[msg] (68,-108) -- ++(5,0) -- ++(0,-3) -- ++(-5,0);
  \node[font=\scriptsize,anchor=west] at (74,-109.5) {forget conversations, snapshots};
  \draw[ret] (68,-117) -- node[label,above]{200 \code{DELETED}} (34,-117);
  \draw[ret] (34,-121) -- node[label,below]{cookie cleared} (0,-121);
\end{tikzpicture}

  \caption{Account deletion. One SQLite transaction covers every row and the audit chain.}\label{fig:deletion}
\end{figure}

\begin{algorithm}[htbp]
  \caption{Erasure (\texttt{SqliteAccountErasure.erase} in \texttt{sqlite.py}).}\label{alg:erase}
  \KwIn{user $u$, keep\_account (true for a demo persona's reset), erasure draft $d$ with no user id}
  \KwOut{rows deleted per table, number of tombstones}
  \textsc{begin immediate}\;
  \ForEach{table in statements, transactions, receipts, connections, connection\_results, savings\_rules, proposals, authorizations, investment\_links, investment\_snapshots, investment\_documents}{delete rows of $u$; count}
  \lIf{not keep\_account}{delete $u$'s sessions, profile and account; count}
  $E \gets \emptyset$\;
  \ForEach{audit event $e$ of $u$, by sequence}{
    \lIf{$e$ is an erasure event}{\KwContinue}
    replace $e$ with its tombstone naming $d$ (same sequence, \code{prev_hash}, \code{event_hash}); add $e.\mathit{sequence}$ to $E$\;
  }
  append $d$ with detail \texttt{erased\_events=$|E|$; erased\_sha256=SHA256($E$); counts}\;
  \textsc{commit}; WAL checkpoint (truncate)\;
\end{algorithm}

\textbf{Demo personas are reset, not deleted.} Camila and Diego are shared fixtures with no password. In a persona
session the profile page offers \emph{Reset demo data}: confirmation \code{RESET}, no password, and the erasure above
with \code{keep_account} true, recorded as \code{DEMO_DATA_RESET}. The fixture and every visitor's session stay.

\subsection{Evidence}

\code{make test-security} runs the attack tests S01--S24, S30--S32 and S35--S38; \code{docs/SECURITY.md} lists each
with its expected behaviour and audit event, except S24 (concept cards), whose row pull request \#58 adds
(\cref{sec:discrepancies}). The newer ones cover
concept-card injection and the \code{EDUCATIONAL_NOT_ADVICE} warning (S24), Fintual goal names, credentials and imported
statements (S30--S32), and connected servers: false read-only labels, injected answers, SSRF and token leaks
(S35--S38). Receipts have their own attack tests in \code{tests/security/test_receipts.py}. The suitability suite
(\code{make test-suitability}) now passes 24 checks with no expected failures, UG01--UG03 included. Option B adds \code{tests/adapters/test_data_key.py} (key sources, deployed refusals, wrong key,
migration), \code{tests/adapters/test_encryption_at_rest.py} (raw database bytes after a full HTTP flow; upgrade of a
schema~0 file) and \code{tests/security/test_account_deletion.py} (user B intact, no row of any table refers to user A,
forged tombstones fail). \code{tests/security/test_compliance.py} turns each legal requirement into a test that names
its norm (\cref{sec:privacy}).
